\documentclass[../../main.tex]{subfiles}
\begin{document}
% SOURCE p=233; printed=218; object=ch07-semicontinuity-definition; region=middle
本节最后一组结果讨论半连续函数。
\begin{definition}[下半连续与上半连续]\label{definition:ch07-semicontinuity}
设 $X$ 为拓扑空间。函数 $f:X\to(-\infty,\infty]$ 称为\textbf{下半连续}（LSC），若对所有 $a\in\mathbb R$，集合 $\{x:f(x)>a\}$ 都开。函数 $f:X\to[-\infty,\infty)$ 称为\textbf{上半连续}（USC），若对所有 $a\in\mathbb R$，集合 $\{x:f(x)<a\}$ 都开。
\end{definition}
% SUPPLEMENT: id=v14-ch07-semicontinuity-definition; parent=ch07-semicontinuity-definition; kind=proof-detail+motivation+hypotheses; origin=independent-derivation
\begin{remark}
\textbf{动机与值域约定（重构新增）。}下半连续允许函数在极限点取较小值，上半连续允许取较大值；例如开集与闭集的示性函数分别体现这两种方向。下半连续函数的严格上水平集开，故它 Borel 可测；上半连续同理，因为开半直线生成相应扩展实直线的 Borel 结构。按本书约定，下半连续可取 $+\infty$ 而不取 $-\infty$，上半连续则相反，因而同一类两个函数求和不会出现 $\infty-\infty$。取负号交换这两种类：$h$ 上半连续当且仅当 $-h$ 下半连续。实值函数连续当且仅当同时上、下半连续，因为开区间的原像是两个上述开原像的交。
\end{remark}
% END SUPPLEMENT: id=v14-ch07-semicontinuity-definition

% SOURCE p=233; printed=218; object=ch07-lsc-properties; number=7.11; region=middle-lower
\begin{proposition}\label{proposition:ch07-lsc-properties}
设 $X$ 为拓扑空间。
\begin{enumerate}[label=\alph*.,ref=(\alph*)]
\item\label{item:ch07-open-indicator-lsc} 若 $U$ 开，则 $\chi_U$ 下半连续。
\item\label{item:ch07-lsc-scaling} 若 $f$ 下半连续且 $c\in[0,\infty)$，则 $cf$ 下半连续。
\item\label{item:ch07-lsc-supremum} 若 $\mathscr G$ 是一族下半连续函数，$f(x)=\sup\{g(x):g\in\mathscr G\}$，则 $f$ 下半连续。
\item\label{item:ch07-lsc-sum} 若 $f_1,f_2$ 下半连续，则 $f_1+f_2$ 下半连续。
\item\label{item:ch07-lsc-continuous-minorants} 若 $X$ 为 LCH 空间且 $f$ 非负下半连续，则
\begin{align*}
f(x)=\sup\{g(x):g\in C_c(X),\ 0\leqslant g\leqslant f\}.
\end{align*}
\end{enumerate}
\end{proposition}
% SOURCE p=233; printed=218; object=ch07-lsc-properties-proof; region=lower
\begin{proof}
(a)、(b) 显然，(c) 来自
\begin{align*}
f^{-1}((a,\infty])=\bigcup_{g\in\mathscr G}g^{-1}((a,\infty]).
\end{align*}
对 (d)，若 $f_1(x_0)+f_2(x_0)>a$，选 $\varepsilon>0$ 使 $f_1(x_0)>a-f_2(x_0)+\varepsilon$，则
\begin{align*}
\{x:(f_1+f_2)(x)>a\}\supseteq{}&\{x:f_1(x)>a-f_2(x_0)+\varepsilon\}\cap{}\\
&\{x:f_2(x)>f_2(x_0)-\varepsilon\},
\end{align*}
右端为 $x_0$ 的邻域。因此 $f_1+f_2$ 下半连续。最后，若 $X$ 为 LCH 空间、$f(x)>0$ 且 $0<a<f(x)$，则 $U=\{y:f(y)>a\}$ 是含 $x$ 的开集。由 Urysohn 引理，存在 $g\in C_c(X)$ 使 $g(x)=a$ 且 $0\leqslant g\leqslant a\chi_U\leqslant f$。令 $a\uparrow f(x)$ 即证 (e)；$f(x)=0$ 时结论平凡。
\end{proof}
% SUPPLEMENT: id=v14-ch07-lsc-properties-details; parent=ch07-lsc-properties; kind=omitted-proof+proof-detail+correction+motivation+hypotheses; origin=independent-derivation
\begin{remark}
\textbf{端点约定与无穷值补证（重构新增）。}(a) 中 $\{\chi_U>a\}$ 在 $a<0$、$0\leqslant a<1$、$a\geqslant1$ 三种情形分别为 $X,U,\varnothing$，均开。(b) 中 $c>0$ 时 $\{cf>a\}=\{f>a/c\}$；$c=0$ 时约定 $0\cdot(+\infty)=0$，所得零函数连续。负系数会交换上、下半连续方向，不能直接保留结论。(c) 按本书值域须理解为\emph{非空}函数族；任取其中一个 $g$，就有 $f(x)\geqslant g(x)>-\infty$。空族上确界通常为 $-\infty$，不在这里规定的下半连续值域内。任意族并集仍开，所以无需族可数。

原 (d) 的阈值写法在 $f_2(x_0)=+\infty$ 时不能直接使用。适用于所有允许值的完整证明是
\begin{align*}
\{f_1+f_2>a\}=\bigcup_{r\in\mathbb Q}\bigl(\{f_1>r\}\cap\{f_2>a-r\}\bigr).
\end{align*}
右到左来自相加；左到右若两值有限，在 $a-f_2(x)<r<f_1(x)$ 中选有理数，若某值为 $+\infty$，同样可选足够大或足够小的有限 $r$；两者都为 $+\infty$ 时任意 $r$ 都可。故该原像开。(e) 的候选族含零函数，所以在 $f(x)=0$ 处取界为零；若 $f(x)=+\infty$，上述任意有限 $a>0$ 的峰函数使上确界无穷。所有候选不超过 $f$，给出另一方向的界。

动机是把开水平集转为紧支连续的下方测试。前四项只用拓扑；(e) 特别需要 LCH 的局部 Urysohn 引理。$f\geqslant0$ 保证零候选和峰函数在 $U^c$ 上也不超过 $f$，不能在原式中随意去掉。函数族的上确界允许无穷，正是许多积分极限的自然值域。
\end{remark}
% END SUPPLEMENT: id=v14-ch07-lsc-properties-details

% SOURCE p=233; printed=218; object=ch07-usc-dual-transition; region=bottom
上半连续函数当然也有相应结果；原书把其陈述留给读者。下一个结果是下半连续函数网的单调收敛定理。
% SUPPLEMENT: id=v14-ch07-usc-dual-properties; parent=ch07-usc-dual-transition; kind=omitted-proof+motivation+hypotheses; origin=independent-derivation
\begin{remark}
\textbf{上半连续的全部对偶结论与证明（重构新增）。}若 $F$ 闭，则 $\chi_F$ 上半连续：其严格下水平集在 $a\leqslant0$、$0<a\leqslant1$、$a>1$ 时分别为 $\varnothing,X\setminus F,X$。若 $h$ 上半连续且 $c\geqslant0$ 有限，则 $ch$ 上半连续；$c>0$ 时用阈值 $a/c$，$c=0$ 时约定 $0\cdot(-\infty)=0$。非空上半连续族的逐点下确界仍上半连续，因为 $\{\inf h<a\}=\bigcup_h\{h<a\}$；非空保证下确界不取 $+\infty$。两个上半连续函数之和仍上半连续，由它们的负函数下半连续及上一命题的加法结论得到。

若 $X$ 为 LCH 空间且 $h\leqslant0$ 上半连续，则
\begin{align*}
h(x)=\inf\{g(x):g\in C_c(X),\ h\leqslant g\leqslant0\}.
\end{align*}
事实上 $-h$ 非负下半连续，应用\subcref{proposition:ch07-lsc-properties}{item:ch07-lsc-continuous-minorants}后取负号，候选 $u$ 与 $g=-u$ 一一对应，便得此式；这也包含 $h(x)=-\infty$ 的情形。此处必须同时反转符号、取界方向和点值约束，不能把非负下逼近的原式只机械替换 LSC 为 USC。闭集示性函数虽非负，仍上半连续；它将作为稍后非负可测函数的有限下逼近，而上述对偶公式讨论的是非正函数的连续上逼近。
\end{remark}
% END SUPPLEMENT: id=v14-ch07-usc-dual-properties

% SOURCE p=234; printed=219; object=ch07-directed-lsc-convergence; number=7.12; region=upper
\begin{proposition}\label{proposition:ch07-directed-lsc-convergence}
设 $\mathscr G$ 是 LCH 空间 $X$ 上按 $\leqslant$ 向上有向的一族非负下半连续函数，即对任意 $g_1,g_2\in\mathscr G$，存在 $g\in\mathscr G$ 使 $g_1\leqslant g$ 且 $g_2\leqslant g$。令 $f=\sup\{g:g\in\mathscr G\}$。若 $\mu$ 为 $X$ 上的任意 Radon 测度，则
\begin{align*}
\int f\,\mathrm d\mu=\sup\left\{\int g\,\mathrm d\mu:g\in\mathscr G\right\}.
\end{align*}
\end{proposition}
% SOURCE p=234; printed=219; object=ch07-directed-lsc-convergence-proof; region=upper-middle
\begin{proof}
由\subcref{proposition:ch07-lsc-properties}{item:ch07-lsc-supremum}，$f$ 下半连续，从而 Borel 可测；显然 $\int f\,\mathrm d\mu\geqslant\sup_{g\in\mathscr G}\int g\,\mathrm d\mu$。为证反向，考虑\cref{theorem:ch02-simple-approximation}证明中构造的、逐点递增到 $f$ 的简单函数
\begin{align*}
\phi_n=2^{-n}\sum_{j=1}^{2^{2n}}\chi_{U_{nj}},\qquad U_{nj}=\{x:f(x)>j2^{-n}\}.
\end{align*}
由单调收敛定理，给定 $\alpha<\int f\,\mathrm d\mu$，可固定充分大的 $n$ 使 $2^{-n}\sum_j\mu(U_{nj})=\int\phi_n\,\mathrm d\mu>\alpha$。各 $U_{nj}$ 开，所以存在紧集 $K_j\subseteq U_{nj}$ 使 $2^{-n}\sum_j\mu(K_j)>\alpha$。令 $\psi=2^{-n}\sum_j\chi_{K_j}$。对每个 $x\in\bigcup_jK_j$，有 $f(x)>\phi_n(x)\geqslant\psi(x)$，故可选 $g_x\in\mathscr G$ 使 $g_x(x)>\psi(x)$。因为 $-\chi_{K_j}$ 下半连续，$g_x-\psi$ 也下半连续，所以 $V_x=\{y:\psi(y)<g_x(y)\}$ 开。这些 $V_x$ 覆盖紧集 $\bigcup_jK_j$，可取有限子覆盖 $V_{x_1},\ldots,V_{x_m}$。选 $g\in\mathscr G$ 使 $g_{x_k}\leqslant g$ 对所有 $k$ 成立，则 $\psi\leqslant g$，故 $\int g\,\mathrm d\mu>\alpha$。由于 $\alpha$ 任意，结论成立。
\end{proof}
% SUPPLEMENT: id=v14-ch07-directed-lsc-details; parent=ch07-directed-lsc-convergence; kind=proof-detail+motivation+hypotheses; origin=independent-derivation
\begin{remark}
\textbf{有向性如何替代序列（重构新增）。}按有向集约定 $\mathscr G\ne\varnothing$；它保证有限子族也有共同上界，证明可对子族大小归纳。若 $\int f=0$，结论由非负性直接成立；其余情况只需处理 $0\leqslant\alpha<\int f$。各 $\phi_n$ 确实递增：在第 $n$ 级命中的阈值 $j2^{-n}$，第 $n+1$ 级对应的 $(2j-1)2^{-n-1}$、$2j2^{-n-1}$ 都命中，二者权重之和为原权重。若 $0<f(x)\leqslant2^n$，命中层数给出 $0<f(x)-\phi_n(x)\leqslant2^{-n}$；若 $f(x)>2^n$，则 $\phi_n(x)=2^n<f(x)$。因此 $\phi_n\uparrow f$，并且 $f>\phi_n$ 在 $\{f>0\}$ 上\emph{严格}成立。这解释为何原构造使用 $>$ 水平集而不是 $\geqslant$。

固定 $n$ 后只有有限个开集，内正则性允许同时选各紧集使其加权和仍大于 $\alpha$：有限总和时把正余量分配给各项；某一项无限时，仅该项就可选足够大的有限紧质量。各 $K_j$ 闭，故 $-\chi_{K_j}$ 下半连续，$-\psi$ 为有限值下半连续函数，所以 $g_x-\psi$ 没有无穷值相减问题。在紧并之外 $\psi=0\leqslant g$，在紧并之内用有限覆盖及共同上界，才得到全空间的比较。

动机是把一个不可数上确界的积分压缩到有限个局部测试。一般可测函数的不可数有向族没有这样的结论：在 $[0,1]$ 上，有限子集的示性函数向上有向，每个积分为零，上确界却恒为一。它们通常不是下半连续。这里下半连续性产生开覆盖，紧内逼近给出有限子覆盖，非负性控制紧集外并保证扩展值积分合法。没有假设 $\mathscr G$ 可数，也不要求 $\mu$ 总质量有限。
\end{remark}
% END SUPPLEMENT: id=v14-ch07-directed-lsc-details

% SOURCE p=234; printed=219; object=ch07-lsc-integral-minorants; number=7.13; region=middle
\begin{corollary}\label{corollary:ch07-lsc-integral-minorants}
若 $\mu$ 为 Radon 测度，$f$ 非负下半连续，则
\begin{align*}
\int f\,\mathrm d\mu=\sup\left\{\int g\,\mathrm d\mu:g\in C_c(X),\ 0\leqslant g\leqslant f\right\}.
\end{align*}
\end{corollary}
% SOURCE p=234; printed=219; object=ch07-lsc-integral-minorants-proof; region=middle
\begin{proof}
结合\subcref{proposition:ch07-lsc-properties}{item:ch07-lsc-continuous-minorants}与\cref{proposition:ch07-directed-lsc-convergence}即可。
\end{proof}
% SUPPLEMENT: id=v14-ch07-lsc-integral-minorants-details; parent=ch07-lsc-integral-minorants; kind=proof-detail+motivation+hypotheses; origin=independent-derivation
\begin{remark}
\textbf{候选族的完整核验（重构新增）。}令 $\mathscr G=\{g\in C_c(X):0\leqslant g\leqslant f\}$。零函数在其中，故非空。若 $g_1,g_2\in\mathscr G$，则 $\max(g_1,g_2)$ 连续，支集包含于两紧支集之并，并仍介于零与 $f$ 之间，故是族内共同上界。每个连续函数都下半连续，而该族的逐点上确界由前命题正好为 $f$。因此有向收敛命题的全部假设成立，给出公式。每个候选自身可积，因为紧集有限；其积分上确界却允许无穷。这个结论把开集的连续内逼近提升到一般非负下半连续函数，是表示定理与积分逼近之间的桥梁。
\end{remark}
% END SUPPLEMENT: id=v14-ch07-lsc-integral-minorants-details

% SOURCE p=234; printed=219; object=ch07-semicontinuous-integral-approximation; number=7.14; region=lower
\begin{proposition}\label{proposition:ch07-semicontinuous-integral-approximation}
若 $\mu$ 为 Radon 测度，$f$ 非负 Borel 可测，则
\begin{align*}
\int f\,\mathrm d\mu=\inf\left\{\int g\,\mathrm d\mu:g\geqslant f,\ g\text{ 下半连续}\right\}.
\end{align*}
若 $\{x:f(x)>0\}$ 为 $\sigma$ 有限集，则还有
\begin{align*}
\int f\,\mathrm d\mu=\sup\left\{\int g\,\mathrm d\mu:0\leqslant g\leqslant f,\ g\text{ 上半连续}\right\}.
\end{align*}
\end{proposition}
% SOURCE p=234; printed=219; object=ch07-semicontinuous-integral-approximation-proof; region=bottom
\begin{proof}
取非负简单函数列 $\phi_n\uparrow f$。由 $f=\phi_1+\sum_{n=2}^\infty(\phi_n-\phi_{n-1})$，每一项均为非负简单函数，可以写成
\begin{align*}
f=\sum_{j=1}^\infty a_j\chi_{E_j},\qquad a_j>0.
\end{align*}
给定 $\varepsilon>0$，为每个 $j$ 取开集 $U_j\supseteq E_j$ 使 $\mu(U_j)\leqslant\mu(E_j)+\varepsilon/(2^ja_j)$。令 $g=\sum_j a_j\chi_{U_j}$，由\cref{proposition:ch07-lsc-properties}，$g$ 下半连续；同时 $g\geqslant f$ 且 $\int g\,\mathrm d\mu\leqslant\int f\,\mathrm d\mu+\varepsilon$，得第一个断言。

对第二个断言，给定 $\alpha<\int f\,\mathrm d\mu$，取充分大的 $N$ 使 $\sum_{j=1}^N a_j\mu(E_j)>\alpha$。各 $E_j$ 为 $\sigma$ 有限集，由\cref{proposition:ch07-sigma-finite-inner}可取紧集 $K_j\subseteq E_j$ 使 $\sum_{j=1}^N a_j\mu(K_j)>\alpha$。于是 $g=\sum_{j=1}^N a_j\chi_{K_j}$ 上半连续、$0\leqslant g\leqslant f$ 且 $\int g\,\mathrm d\mu>\alpha$，得结论。
\end{proof}
% SUPPLEMENT: id=v14-ch07-semicontinuous-approximation-details; parent=ch07-semicontinuous-integral-approximation; kind=proof-detail+motivation+hypotheses; origin=independent-derivation
\begin{remark}
\textbf{上下逼近的非对称性（重构新增）。}分解中的 $E_j$ 都可取在 $\{f>0\}$ 内；把每个简单增量的非零水平集枚举即得，有限分解可用空集补齐，零函数也可单独处理。$E_j$ 不必互不相交；逐点非负级数和及单调收敛保证 $\int f=\sum_j a_j\mu(E_j)$。

上逼近若 $\int f=\infty$，任何 $g\geqslant f$ 的积分都无穷，且候选 $g\equiv+\infty$ 存在，所以等式直接成立。若 $\int f<\infty$，每个 $\mu(E_j)<\infty$，外正则可选所述 $U_j$；级数的有限部分和下半连续，非负性使其递增，上确界仍下半连续，且
\begin{align*}
\int g\,\mathrm d\mu=\sum_j a_j\mu(U_j)
\leqslant\sum_j a_j\mu(E_j)+\varepsilon\sum_j2^{-j}
=\int f\,\mathrm d\mu+\varepsilon.
\end{align*}
下逼近只需有限部分和，有限个紧集的加权示性函数为\emph{有限值}上半连续函数；候选零函数处理 $\int f=0$，其余情况取 $0\leqslant\alpha<\int f$ 即可。若有限部分和中某项积分无穷，其内正则性允许单独取足够大的紧质量；若各项有限，则分配有限的正余量。于是无限积分情形也被涵盖。

动机是把集合的开外逼近和紧内逼近分别提升成函数的下半连续上包及上半连续下包。外正则对全部 Borel 集成立，所以上式不加 $\sigma$ 有限条件；内正则只已知对 $\sigma$ 有限 Borel 集成立，所以下式需该条件。有限积分自动使 $\{f>0\}=\bigcup_n\{f\geqslant1/n\}$ 为 $\sigma$ 有限，因为 $\mu\{f\geqslant1/n\}\leqslant n\int f$。因此可积非负函数总可使用两侧逼近。
\end{remark}
% END SUPPLEMENT: id=v14-ch07-semicontinuous-approximation-details
\subfile{section-02-exercises.tex}
\end{document}
